\documentclass[11pt]{article}
\usepackage[a4paper,margin=1in]{geometry}
\usepackage{amsmath,amssymb,mathtools,bm}
\usepackage{enumitem,booktabs,array}
\usepackage{tcolorbox}
\usepackage{hyperref}
\usepackage[protrusion=true,expansion=false]{microtype}
\hypersetup{colorlinks=true,linkcolor=blue,urlcolor=blue,citecolor=blue}
\setlist[itemize]{topsep=2pt,itemsep=2pt}
\setlist[enumerate]{topsep=2pt,itemsep=2pt}
\newtcolorbox{keybox}{colback=blue!3!white,colframe=blue!40!black,boxrule=0.4pt,arc=1mm}
\newtcolorbox{alertbox}{colback=red!3!white,colframe=red!40!black,boxrule=0.4pt,arc=1mm}
\newtcolorbox{answerbox}{colback=green!3!white,colframe=green!40!black,boxrule=0.4pt,arc=1mm}
\newtcolorbox{drillbox}{colback=yellow!5!white,colframe=orange!60!black,boxrule=0.4pt,arc=1mm}
\newcommand{\EV}{\mathbb{E}}
\newcommand{\Var}{\mathrm{Var}}
\newcommand{\Cov}{\mathrm{Cov}}
\newcommand{\blank}[1]{\underline{\hspace{#1}}}

\title{E11-03: Farre-Mensa \& Ljungqvist (2016, RFS) \\
\large ``Do Measures of Financial Constraints Measure Financial Constraints?'' --- Active Recall Workbook}
\author{Empirical Corporate Finance II --- Exam Prep}
\date{\today}
\begin{document}\maketitle

\begin{keybox}
\textbf{Paper in one sentence.} Using quasi-experimental state-tax changes (as in Heider-Ljungqvist), FML test whether the four most-popular financial-constraints indices (Kaplan-Zingales, Whited-Wu, Hadlock-Pierce, payout-ratio) identify firms whose investment responds to cash-flow shocks. Answer: none of the indices reliably separate constrained from unconstrained firms --- they all pick up firm attributes correlated with growth opportunities, not constraint status.

\textbf{Placement.} Landmark methodological paper challenging the empirical literature that relies on financial-constraint indices; paradigm-shifting for cash-flow sensitivity research.
\end{keybox}

\section*{Round 1: Complete Walk-Through}

\subsection*{1.1 Motivation}
Financial-constraints research (Fazzari-Hubbard-Petersen 1988, Kaplan-Zingales 1997) classifies firms as constrained/unconstrained using indices and studies differential investment-cash-flow sensitivity. If these indices do not actually measure constraint status, the entire literature is compromised.

\subsection*{1.2 Test idea}
If a firm is truly financially constrained, a positive (exogenous) cash-flow shock should raise its investment, whereas an unconstrained firm should invest at its first-best level regardless. Use state-tax cuts as exogenous cash-flow shocks and ask which groups (by various indices) respond.

\subsection*{1.3 Data and design}
\begin{itemize}
\item US state corporate-income-tax changes 1989--2011.
\item Firms classified by KZ, WW, HP, payout-ratio indices.
\item Event-study of investment response to tax-induced cash-flow changes across constraint groups.
\end{itemize}

\subsection*{1.4 Headline results}
\begin{itemize}
\item Investment response to cash-flow shocks is similar across all constraint-index quartiles.
\item Firms classified as ``most constrained'' by KZ or WW do not respond more strongly than ``least constrained'' firms.
\item Instead, size (HP index) and growth opportunities predict investment response more consistently.
\item Payout-ratio index does a slightly better job but still leaves huge variance unexplained.
\end{itemize}

\subsection*{1.5 Interpretation}
\begin{itemize}
\item Popular constraint indices fail to identify financially-constrained firms as defined by investment-cash-flow sensitivity.
\item Indices pick up attributes (growth, size, age) correlated with but distinct from financial constraints.
\item Empirical-corporate-finance studies using these indices need caution.
\end{itemize}

\subsection*{1.6 Implications}
\begin{itemize}
\item Use exogenous cash-flow shocks (IV or quasi-experiment) rather than indices for constraint research.
\item Re-examine prior results that relied on KZ/WW/HP classifications.
\item Develop better direct measures (e.g., covenant violations, rating downgrades, rationing in credit markets).
\end{itemize}

\subsection*{1.7 Caveats}
\begin{alertbox}
\begin{itemize}
\item State-tax shocks measure short-run cash flow, not long-run constraint.
\item Sample is US-listed; private firms may differ.
\item Indices may still correlate with constraints for specific use cases (e.g., firm survival).
\end{itemize}
\end{alertbox}

\section*{Round 2: Level 1 Fill-in}
\begin{enumerate}
\item Four tested indices: KZ, WW, \blank{1cm}, payout-ratio.
\item Shock: US \blank{2cm} corporate-tax changes.
\item Finding: indices do \blank{2cm} identify constraint status.
\item Implication: use exogenous \blank{2cm} for constraint research.
\item HP stands for \blank{2cm} (authors' initials).
\end{enumerate}

\section*{Round 3: Level 2 Prompts}
\begin{enumerate}
\item Describe the identification idea (tax-induced cash-flow shock as test).
\item Which indices are tested and how they were constructed.
\item Report main findings on differential response.
\item Discuss implications for financial-constraints literature.
\item What alternative measures should researchers use?
\end{enumerate}

\section*{Round 4: Minimal Scaffolding}
From a blank page: (i) motivation, (ii) test, (iii) findings, (iv) implications, (v) alternatives.

\section*{Round 5: Blank Page}
Essay: measuring financial constraints --- a methodological critique.

\vspace{6pt}
\begin{drillbox}
\textbf{Speed Drill.} (1) Indices tested. (2) Shock. (3) Finding. (4) Bias mechanism. (5) Alternative measure.
\end{drillbox}

\section*{Quick Reference \& Answer Key}
\begin{answerbox}
\begin{itemize}
\item \textbf{Indices.} KZ, WW, HP, payout ratio.
\item \textbf{Shock.} State tax changes as exogenous CF shocks.
\item \textbf{Finding.} Indices fail to identify constraint status.
\item \textbf{Advice.} Use exogenous shocks, not indices.
\end{itemize}
\end{answerbox}

\begin{keybox}
\textbf{30-second exam pitch.} Farre-Mensa \& Ljungqvist (2016) test whether the four leading financial-constraints indices (Kaplan-Zingales, Whited-Wu, Hadlock-Pierce, payout ratio) actually identify financially-constrained firms. Using US state corporate-tax changes 1989--2011 as exogenous cash-flow shocks, they compare the investment response across index quartiles. Under the classical constraint hypothesis, ``constrained'' firms should show stronger investment responses. Instead, response is similar across all quartiles, and the indices primarily capture growth opportunities, size, or age rather than constraint status. The paper is a paradigm-shifting critique of decades of empirical work using these indices, pushing researchers toward quasi-experimental cash-flow shocks and direct measures (covenant violations, credit rationing) rather than index-based classification.
\end{keybox}

\end{document}
